% Source: https://artofproblemsolving.com/community/c3h2945150_san_jose_math_circle_monthly_contest_201819_intermediate_contest_34
% Stable identity: aops-topic-2945150-q02
First prove $A\ge14$. If two consecutive rows contain no removed cell, their full-width rectangle has area 14. Otherwise the only way at most three bad rows can separate all four or more good rows is for the bad rows to be exactly 2,4,6. Each then has exactly one removed cell. If any removed cell has column at most 2 or at least 6, the three-row strip centered on its row has no other removed cell and leaves five consecutive intact columns on one side, giving area at least 15. Otherwise all three removed cells lie in columns 3,4,5, so the first two columns are intact through all seven rows, giving 14. For sharpness remove cells at (column,row) $(3,2),(5,4),(3,6)$. For rectangle heights 1 through 7, the largest achievable areas are respectively 7,8,12,8,10,12,14: for each consecutive row interval, exclude its removed columns and use its longest remaining consecutive column block. The following 28-interval check certifies that table. Thus $\boxed{A=14}$.
\begin{verbatim}
holes=[(2,1),(4,3),(2,5)] # zero-based (column,row)
best={h:0 for h in range(1,8)}
for lo in range(7):
 for hi in range(lo+1,8):
  blocked=sorted({x for x,y in holes if lo<=y<hi})
  width=max(b-a-1 for a,b in zip([-1]+blocked,blocked+[7]))
  best[hi-lo]=max(best[hi-lo],width*(hi-lo))
assert list(best.values())==[7,8,12,8,10,12,14]
\end{verbatim}

\begin{asy}
usepackage("amsmath");
size(7cm);
for(int x=0;x<7;++x)for(int y=0;y<7;++y){path p=(x,y)--(x+1,y)--(x+1,y+1)--(x,y+1)--cycle;bool bad=(x==2&&(y==1||y==5))||(x==4&&y==3);filldraw(p,bad?gray:white,black);}draw((0,0)--(2,0)--(2,7)--(0,7)--cycle,blue+1bp);label("an intact area14 rectangle",(3.5,-.7));
\end{asy}
